\documentclass[11pt]{article}
\input{../suite_preamble}
\input{../suite_macros}

\title{Reflective Non-Exhaustion:\\[0.35em]
\large Certification, Realization, Fibers, and Obstruction}
\author{Nova Spivack}
\date{2026}

\begin{document}

\maketitle

\begin{abstract}
We state and defend a summit-level theoremic claim emerging from a machine-checked
program in Lean~4: \emph{canonical certification does not, in general, exhaust
realization}. More precisely, when a mathematical or formal architecture contains
a bare certification layer, an enriched realization layer, and a sound comparison
map from the latter to the former, collapse at the bare carrier may coexist with
nontrivial realized structure above that collapse. When this occurs, the residue
is not amorphous: it is organized by non-injective forgetful comparison, nontrivial
fibers, sections, and obstruction laws.

This paper does not introduce new axioms or replace the companion formal papers.
Its function is different. It distills the permanent idea of the program into a
single theoremic synthesis, states exactly what has been proved across the suite,
explains why the result is not merely a restatement of G\"odelian incompleteness or
a generic observation about proof relevance in dependent type theory, and argues
that the collapse/refinement/fiber/obstruction architecture constitutes a reusable
research instrument. The claim is supported by five completed kinds of discharge:
the internal origin theorem in Infinity Compression, external transfer across
twelve benchmark families, an algebraic obstruction theorem for group extensions,
an arithmetic bridge through embedding problems, a topological discharge through
Quillen's Theorem~A for Galois connections, and a computability anchor via halting
and Rice-style non-certifiability results. We conclude with the exact scope of the
claim and the frontier required to sharpen it into its strongest universal form.
\end{abstract}

\SuiteSeriesContext{This manuscript is the \emph{summit statement}.}

\keywords{formal verification, certification, realization, fibers, obstruction,
  reflective systems, Lean~4, Mathlib}
\amsclassification{03B70, 18A30, 18G50, 68V15, 68V20}

\tableofcontents
\newpage

\section{The summit claim}\label{sec:summit-claim}

\subsection{The central statement}

The program distilled in this paper is centered on the following claim.

\begin{theorem}[Reflective Non-Exhaustion, program statement]\label{thm:rne}
Let a \emph{reflective certification architecture} consist of:
\begin{enumerate}[label=(\roman*)]
  \item a bare certification carrier $B$;
  \item an enriched realization carrier $E$;
  \item a sound comparison map $\pi : E \to B$;
  \item a distinguished canonical region $C \subseteq B$ singled out by the theory's
    certification rules.
\end{enumerate}
Then the following phenomenon can occur, and is realized in the companion formal
developments:
\begin{enumerate}[label=(\alph*)]
  \item the certification layer collapses canonically inside $C$;
  \item the realization layer remains nontrivial above that collapse;
  \item the nontriviality is organized by non-injectivity of $\pi$, by nontrivial
    fibers $\pi^{-1}(b)$ for $b \in C$, and by explicit obstruction or section laws
    governing when collapse can or cannot be lifted.
\end{enumerate}
In this sense, canonical certification does not exhaust realization.
\end{theorem}

\noindent
The theorem above is a summit-level synthesis statement. Its fully formal
instances are proved in the companion papers and Lean modules cited below; this
paper does not ask the reader to accept it as a slogan detached from formal
mathematics.

\subsection{Public paraphrase}

A concise public rendering is:

\begin{quote}
\emph{A system's canonical certificate need not fully contain the structure by
which that certificate is realized.}
\end{quote}

A stronger but still careful paraphrase is:

\begin{quote}
\emph{Canonical certification does not, in general, exhaust realized structure;
the residue lives in fibers above certification.}
\end{quote}

\subsection{What this paper is for}

This paper is not the place where the detailed Lean proofs live; nor is it a
popular essay. Its role is to do four things that none of the companion papers,
taken alone, can do as sharply:

\begin{enumerate}[label=(\arabic*)]
  \item state the permanent theoremic idea in one place;
  \item show exactly which completed results already support it;
  \item separate it from nearby but different ideas (incompleteness, generic proof
    relevance, informal rhetoric);
  \item explain why the architecture is a reusable research instrument rather than
    a one-off observation.
\end{enumerate}

\section{What is actually proved in the suite}\label{sec:proved-suite}

The summit claim is supported by several completed theorem families, all formalized
in Lean~4 with zero \texttt{sorry} in the relevant modules.

\subsection{Internal origin theorem}

The flagship Infinity Compression paper~\cite{SpivackIC} proves the core
certification/realization split on named carriers internal to the library:

\begin{enumerate}[label=(\arabic*)]
  \item canonical collapse at the bare certification carrier;
  \item a reflective split theorem on enriched witnesses;
  \item strict refinement via a typed forgetful map;
  \item nontriviality of the canonical fiber;
  \item injective role coordinates on a standard fiber patch.
\end{enumerate}

This is the \emph{internal origin theorem}: the place where the architecture first
appears not as vocabulary but as a proved, load-bearing structure.

\subsection{External transfer}

The validation paper~\cite{SpivackValidation} proves that the same packaging
exports to twelve independent mathematical benchmark families: quotients,
products, dependent sums, coproduct discriminants, subtypes, orbit quotients,
ideal quotients, pushouts, classifiers, localization, pullbacks, and a
non-surjective control case.

The role of that paper is not theorem-level domain leverage. Its role is
transferability: it demonstrates that the architecture is not an idiosyncratic
artifact of one flagship library.

\subsection{Algebraic discharge}

The group extension paper~\cite{SpivackFiber} proves the central algebraic
obstruction law:
\[
\text{splits} \iff \text{trivial cocycle}.
\]
It also proves the section cocycle identity, local-to-global obstruction
principles, and associated fiber structure.

This is a genuine theorem-level discharge in classical mathematics: the
architecture is not merely organizational there, but mathematically load-bearing.

\subsection{Arithmetic discharge}

The same paper and associated Lean modules show that embedding problems are
equivalent to extension splitting problems and hence to cocycle triviality:
\[
\text{solvable embedding problem}
\iff
\text{splitting}
\iff
\text{trivial cocycle}.
\]
This supplies the arithmetic face of the program, connecting the architecture to
inverse Galois-style lifting problems.

\subsection{Topological discharge}

The Quillen paper~\cite{SpivackQuillen} proves Quillen's Theorem~A for Galois
connections at the simplicial/nerve level, using contractible fibers and
homotopy-equivalence data. This is the topological instance of the same
architecture: fibers, explicit comparison maps, and a theorem-level consequence.

\subsection{Computability anchor}

The summit-supporting Lean modules also contain a logic/computability anchor:
halting is not capturable by the relevant certification layer, and Rice-style
phenomena block certification of nontrivial semantic properties in the
\texttt{ComputablePred} setting. This does not mean the program is ``just
G\"odel again,'' but it ensures the summit architecture touches standard
non-certifiability barriers in classical computability theory.

\subsection{Summary of support}

Taken together, the suite already contains:
\begin{center}
\begin{tabular}{@{}ll@{}}
\toprule
\textbf{Support type} & \textbf{Completed instance} \\
\midrule
Internal origin & certification vs.\ realization split in IC \\
Transferability & twelve external benchmark families \\
Algebraic discharge & group extension splitting criterion \\
Arithmetic discharge & embedding problems and tower bridge \\
Topological discharge & Quillen A for Galois connections \\
Computability anchor & halting / Rice non-certifiability layer \\
\bottomrule
\end{tabular}
\end{center}

That is enough to justify a summit paper. The summit claim is no longer a mere
program aspiration.

\section{The architecture itself}\label{sec:architecture}

\subsection{Certification}

A \emph{certification layer} is a carrier of theorem-level or officially packaged
invariants. It is the layer at which a theory says: this is the certified object,
certificate, derivation, witness, or verdict.

The defining feature of certification, for the present program, is not simply that
it exists, but that it may \emph{collapse}: many legitimate routes may project to
the same canonical certification object.

\subsection{Realization}

A \emph{realization layer} is an enriched carrier above the certificate. It tracks
the structure by which the certificate is obtained or instantiated: sections,
role assignments, dependency shapes, enriched reflective witnesses, splitting
data, chosen representatives, homotopy witnesses, and similar objects.

Realization is not decoration. It becomes mathematically serious exactly when the
comparison map from realization to certification is not injective.

\subsection{Comparison map}

The \emph{comparison map} or \emph{forgetful map}
\[
\pi : E \to B
\]
forgets enriched structure and records only the bare certified image.

This map is the central organizing object. Once $\pi$ is explicit, one can ask:
\begin{enumerate}[label=(\alph*)]
  \item Is $\pi$ surjective?
  \item Is $\pi$ injective?
  \item What do the fibers $\pi^{-1}(b)$ look like?
  \item Are there canonical sections or right inverses?
  \item What obstruction governs lifting from the bare layer back upstairs?
\end{enumerate}

\subsection{Fibers}

The fiber over $b \in B$ is
\[
\Fib(b) := \pi^{-1}(b).
\]
When certification collapses, fibers become the correct place to locate residual
structure. If $\Fib(b)$ has more than one point, then bare certification fails
to determine realization uniquely.

\subsection{Obstruction}

An \emph{obstruction law} is the theorem governing when one may lift from bare to
enriched structure. In different domains this obstruction appears in different
mathematical clothing:
\begin{itemize}
  \item a cocycle in group extensions,
  \item fiber contractibility in Quillen-style comparison,
  \item existence or failure of sections and right inverses in benchmark families,
  \item non-certifiability in computability.
\end{itemize}

The summit claim is that these are not unrelated accidents. They are instances of
a common collapse/refinement/fiber/obstruction architecture.

\section{The internal origin theorem}\label{sec:internal-origin}

The flagship result~\cite{SpivackIC} is the place where the summit architecture is
first proved, rather than merely observed.

\subsection{What happens internally}

At the bare carrier, the standard Phase-2 extraction collapses to a unique canonical
certificate. At the enriched carrier, distinct reflective witnesses remain distinct.
A typed forgetful map from enriched reflective splits to canonical bare certificates
is sound but not injective. The canonical fiber above the bare basepoint is
nontrivial, and role data injects along a standard subfamily.

\subsection{Why that matters}

This gives the program its sharpest internal lesson:

\begin{quote}
Once collapse at the bare carrier is proved, comparison questions that still matter
must move upstairs to enriched carriers and fibers.
\end{quote}

That is a mathematical, not stylistic, conclusion. The internal origin theorem is
what turns the rest of the suite into a coherent program rather than a collection
of analogies.

\section{Transferability}\label{sec:transferability}

The benchmark survey~\cite{SpivackValidation} proves that the architecture exports
to many ordinary mathematical constructions.

\subsection{Why transfer matters}

Without transfer, the flagship theorem could be dismissed as an artifact of an
idiosyncratic internal formalism. The tranche survey blocks that dismissal. It
shows that the same language of:
\[
\text{bare layer} \quad \text{vs.} \quad \text{enriched layer} \quad \text{vs.} \quad
\text{fiber / section / right inverse}
\]
is not local to one formal development.

\subsection{What transfer does not claim}

Transferability is not itself the summit theorem. Nor does it prove new deep
classical results in each benchmark family. Its function is narrower and
essential: it validates the architecture as a reusable interface.

\section{Algebraic and arithmetic discharge}\label{sec:alg-arith}

\subsection{Group extensions}

The group extension theorem
\[
\text{splits} \iff \text{trivial cocycle}
\]
is a model example of the summit architecture:
\begin{itemize}
  \item the bare layer is set-theoretic section existence, which always holds;
  \item the enriched layer is homomorphic splitting;
  \item the comparison is the passage from structured splitting data to bare section data;
  \item the obstruction is the section 2-cocycle.
\end{itemize}

This is exactly the sort of result that makes the architecture mathematically
load-bearing rather than merely descriptive.

\subsection{Embedding problems}

Embedding problems sharpen the picture:
\[
\text{solvable} \iff \text{splits} \iff \text{trivial cocycle}.
\]
Here the same architecture governs an arithmetic lifting problem. Bare data does
not decide liftability by itself; the residue sits in the obstruction.

\subsection{Arithmetic significance}

This is important for the summit claim because it shows that the architecture has
teeth in a domain that is not just internal formal theory and not just generic
category packaging. It reaches into a recognized arithmetic problem class.

\section{Topological discharge}\label{sec:topology}

The Quillen paper~\cite{SpivackQuillen} provides the topological discharge.

\subsection{What happens there}

For Galois connections, the relevant fibers are principal down-sets with terminal
or maximal structure, hence contractible nerves. This gives the comparison theorem:
the induced nerve maps carry explicit homotopy-equivalence data.

\subsection{Architectural reading}

In that setting:
\begin{itemize}
  \item the comparison map is the induced nerve map;
  \item the fibers encode the local comparison geometry;
  \item contractibility is the obstruction-theoretic condition that trivializes
    the comparison problem.
\end{itemize}

This is a very different mathematical domain from group extensions, yet the same
architecture reappears.

\section{Computability anchor}\label{sec:computability}

The summit architecture must not be mistaken for a vague philosophical analogue of
G\"odel. The computability anchor matters because it gives disciplined contact with
classical non-certifiability.

\subsection{What is anchored}

The current Lean support includes halting and Rice-style facts in a certifiability
layer relevant to \texttt{ComputablePred}. These show that there are semantic
phenomena that cannot be fully captured by the relevant certification apparatus.

\subsection{What is not claimed}

The summit theorem is not identical with G\"odel's incompleteness theorem, nor is
the present program claiming to have subsumed all of classical recursion theory.
The role of the computability anchor is narrower: it demonstrates that the program's
language touches standard impossibility terrain in an honest, formal way.

\section{Why this is not just G\"odel}\label{sec:not-godel}

This distinction is central.

\subsection{G\"odel's contribution}

G\"odel shows that sufficiently rich formal systems face unavoidable limits:
undecidable statements, incompleteness, and non-recursive truth phenomena.

\subsection{This program's contribution}

The present program shows something structurally different:
\begin{enumerate}[label=(\arabic*)]
  \item where collapse occurs;
  \item where it provably does not occur;
  \item how the non-collapse is organized.
\end{enumerate}

That is a positive architecture theorem, not merely a barrier theorem.

\subsection{The sharp difference}

G\"odel's theorem says, roughly, that certain truths outrun formal provability.
The summit architecture says that even when a canonical certification exists and
collapses perfectly at one carrier, realized structure may remain irreducibly
nontrivial above it---and that this nontriviality is organized by explicit fibers
and obstructions.

The two ideas are related in spirit, especially once Route~D is included, but they
are not the same theorem.

\section{Why this is not merely proof relevance in type theory}\label{sec:not-dtt}

A common dismissal would be: of course enriched objects exist in dependent type
theory; proof relevance is built in.

That is not the result.

\subsection{What would be trivial}

It is trivial that richer carriers can be defined in DTT. It is trivial that one
can invent a forgetful map from a rich type to a poor one.

\subsection{What is nontrivial here}

What is nontrivial is that the library proves:
\begin{enumerate}[label=(\arabic*)]
  \item collapse at one named carrier;
  \item non-collapse at another named carrier;
  \item non-injectivity of the explicit comparison map;
  \item nontriviality and partial classification of the corresponding fibers;
  \item parallel obstruction laws in external mathematics.
\end{enumerate}

The summit claim is therefore not ``proof relevance exists.'' It is:
\begin{quote}
\emph{this is the correct theoremic architecture for a broad and growing class of
comparison problems.}
\end{quote}

\section{Why this changes mathematical practice}\label{sec:practice}

The strongest reason to take the summit claim seriously is methodological.

\subsection{The reusable recipe}

The architecture gives a reusable research workflow:

\begin{enumerate}[label=(\arabic*)]
  \item identify the bare certificate carrier $B$;
  \item identify an enriched carrier $E$;
  \item define the comparison map $\pi : E \to B$;
  \item prove whether collapse occurs at $B$;
  \item test injectivity or non-injectivity of $\pi$;
  \item if non-injective, study the fibers $\Fib(b)$;
  \item identify sections, right inverses, or obstruction laws governing liftability.
\end{enumerate}

\subsection{Why this matters}

This turns vague statements like ``there is more structure upstairs'' into a
precise theorem program. It tells researchers exactly what to look for and where
to formalize it.

\subsection{Research-instrument claim}

That is why the architecture is best understood not only as a collection of
results, but as a \emph{research instrument}. It gives a stable way to formulate,
compare, and attack problems where theorem-level invariants do not obviously
exhaust the structure that realizes them.

\section{Cross-domain dictionary}\label{sec:dictionary}

A compact dictionary is helpful.

\begin{center}
\small
\begin{tabular}{@{}p{0.18\textwidth}p{0.20\textwidth}p{0.20\textwidth}p{0.18\textwidth}p{0.18\textwidth}@{}}
\toprule
\textbf{Domain} & \textbf{Bare layer} & \textbf{Enriched layer} & \textbf{Comparison map} & \textbf{Fiber / obstruction} \\
\midrule
IC flagship & canonical bare certification & enriched reflective split & typed forgetful map $\pi_A$ & nontrivial canonical fiber \\
Benchmarks & baseline representability & structured witness bundle & projection / forgetting & fibers, sections, right inverses \\
Group extensions & set-theoretic section layer & splitting data & forgetting homomorphic structure & cocycle obstruction \\
Embedding problems & bare target quotient data & lifting / solution data & solution-to-extension comparison & solvable iff trivial cocycle \\
Quillen GC & order-complex / nerve image & homotopy witness data & induced nerve maps & contractible fibers \\
Computability & certifiable predicate layer & realized semantic behavior & certification comparison & halting / Rice barriers \\
\bottomrule
\end{tabular}
\end{center}

This table is not itself a formal theorem. It is the best concise representation
of why the summit architecture is plausible as a unifying language.

\section{Exact scope}\label{sec:scope}

The summit claim is strong, but it is not unlimited.

\subsection{What is claimed}

We claim that:
\begin{enumerate}[label=(\arabic*)]
  \item the collapse/refinement/fiber/obstruction architecture is formally realized
    in the completed companion work;
  \item it transfers across multiple external families;
  \item it yields theorem-level leverage in algebra, arithmetic, topology, and
    computability-adjacent settings;
  \item it is therefore a plausible permanent language for a broad class of
    comparison problems.
\end{enumerate}

\subsection{What is not claimed}

We do \emph{not} claim:
\begin{enumerate}[label=(\arabic*)]
  \item a single current Lean theorem quantifying over all domains at once;
  \item that every formal system must instantiate this exact architecture;
  \item that the current program has already achieved universal adoption or full
    G\"odel-level social impact.
\end{enumerate}

\paragraph{Outside scope.}
Possible speculative applications to self-modeling or phenomenological systems lie outside
the formal scope of this paper; they are not part of the proved theorem content.

\subsection{Why this is enough}

These non-claims do not weaken the summit paper's purpose. A summit statement
should mark the strongest theoremic language already earned by the completed
program, not pretend to have finished all mathematically natural extensions.

\section{Open frontier}\label{sec:frontier}

The next frontier is clear.

\subsection{Mathematical extensions}

The strongest natural next steps are:
\begin{enumerate}[label=(\arabic*)]
  \item fuller universal quantification over a broad abstract architecture class;
  \item stronger fiber-classification theorems;
  \item completion of the $H^1/H^2$ bridge in Mathlib;
  \item deeper arithmetic realization through concrete Galois tower infrastructure;
  \item general Quillen Theorem~A beyond the Galois-connection case.
\end{enumerate}

\subsection{Program-level extensions}

At the level of program closure, the remaining work is not to invent an entirely
new idea, but to strengthen and universalize what is already visible:
\begin{itemize}
  \item sharper inevitability theorems,
  \item stronger unified statements,
  \item broader accepted examples,
  \item publication and adoption.
\end{itemize}

\section{Conclusion}\label{sec:conclusion}

The permanent idea of the program can now be stated cleanly.

Canonical certification may collapse completely at a bare carrier. That collapse
does not guarantee exhaustion of realized structure. When realization remains
nontrivial, the residue is not mysterious: it is organized by explicit comparison
maps, nontrivial fibers, sections, and obstruction laws.

That is the summit claim.

It is not merely a slogan about proof objects. It is not simply a rebranding of
G\"odelian incompleteness. It is not an artifact of one internal library. It is
supported by an internal origin theorem, a transferability survey, theorem-level
discharges in algebra and topology, an arithmetic bridge, and a computability
anchor.

The strongest current reading is therefore:

\begin{quote}
\emph{Collapse downstairs, structured comparison upstairs, and obstruction in the
fibers is a real, reusable theorem architecture.}
\end{quote}

That is why the program deserves a summit paper.

\bibliographystyle{alpha}
\begin{thebibliography}{99}

\bibitem{SpivackIC}
N.~Spivack.
\newblock \emph{Canonical Certification Does Not Exhaust Reflective Structure:
  Reflective Split, Strict Refinement, and Fiber Structure in Infinity Compression}.
\newblock Flagship manuscript, 2026.

\bibitem{SpivackValidation}
N.~Spivack.
\newblock \emph{External Validation of a Positive-Closure Proof Architecture:
  A Multi-Tranche Survey}.
\newblock Validation manuscript, 2026.

\bibitem{SpivackFiber}
N.~Spivack.
\newblock \emph{Fiber Architecture for Group Extensions in Lean~4: Cocycles, Splitting,
  and the Cohomological Bridge}.
\newblock General-method manuscript, 2026.

\bibitem{SpivackQuillen}
N.~Spivack.
\newblock \emph{Quillen's Theorem~A for Galois Connections: A Machine-Checked
  Formalization in Lean~4}.
\newblock General-method manuscript, 2026.

\end{thebibliography}

\end{document}